% SPDX-FileCopyrightText: Copyright (c) 2024-2026 Yegor Bugayenko
% SPDX-License-Identifier: MIT

\section{I/O Streams}\label{sec:streams}

All objects presented in this Section belong to the \ff{io} package.

It is expected that an input stream has the following interface:

\begin{ffcode}
[] > input
  [size] > read /bytes
\end{ffcode}

It is expected that an output stream has the following interface:

\begin{ffcode}
[] > output
  [bytes] > write /true
\end{ffcode}

The \deff{console} object is a basic I/O object that allows interaction
  with the operating system console.
It behaves as ``input'' and ``output.''
%
The \deff{stdout} object is an ``output'' that prints a \deff{string} to the standard output stream.
%
The \deff{stdin} object is an ``input'' and an abstraction of a \deff{string}
  currently available in the standard input stream.
The attribute \deff{next-line} reads a line until the EOL character (\ff{\char`\\n}).
The attribute \deff{all-lines} reads all lines as long as possible.
If there is no string in the stream, the object blocks dataization and waits.
This code endlessly reads strings from the console and immediately prints them back:

\begin{ffcode}
while
  true > [i]
  io.stdout > [i]
    io.stdin.next-line
\end{ffcode}

The object \deff{bytes-as-input} is an ``input'' created from \deff{bytes}.
Reading from it returns the requested number of bytes taken from the wrapped block.
For example, this code wraps a string and reads its first five bytes,
  which decode back to \ff{"Hello"}:

\begin{ffcode}
(io.bytes-as-input "Hello, world!".as-bytes).read 5
\end{ffcode}

The object \deff{malloc-as-output} is an ``output'' directed towards a copy of \deff{malloc}.
Bytes written to it are stored in the memory block, from where they may be read back.
For example, this code writes a string to the output and then prints
  the content of the block as \ff{"Hello!"}:

\begin{ffcode}
malloc.of
  6
  seq * > [m]
    (io.malloc-as-output m).write
      "Hello!".as-bytes
    io.stdout
      string m
\end{ffcode}

The object \deff{tee-input} is an ``input'' and a channel between an ``input''
  and an ``output,'' which moves all available bytes from the former to the
  latter when being dataized.
%
For example, this code writes text to a temporary file:

\begin{ffcode}
(fs.file "/tmp/foo.txt").open
  io.tee-input > [f]
    io.bytes-as-input
      "Hello, world!".as-bytes
    f
\end{ffcode}

The object \deff{input-length} is an object which reads all the bytes
  from the provided ``input'' and returns its length.
%
The object \deff{input-as-bytes} reads all the bytes from the provided ``input''
  and returns them as \deff{bytes}.
%
The object \deff{copied} copies all the bytes from an ``input'' to an ``output''
  and returns the number of bytes copied.
%
The objects \deff{dead-input} and \deff{dead-output} are stubs which read from and write to nothing.
